\documentclass[12pt]{article}
\usepackage[utf8]{inputenc}
\usepackage[T1]{fontenc}
\usepackage{xcolor}



\title{Taller: Librería Universal}
\author{}
\date{}

\begin{document}

\maketitle

La Librería Universal, una entidad de renombre internacional dedicada a la promoción de
obras literarias y académicas de alto impacto, se enfrenta al reto de modernizar su sistema
de catalogación y recomendación. Se desea implementar un sistema digital que no solo
gestione su extensa colección, sino que también ofrezca a sus socios y al público
especializado una herramienta de descubrimiento ágil y precisa.

La información que la Librería considera crucial para cada Libro en su acervo es:

\begin{itemize}
    \item \textbf{Título:} El nombre oficial de la obra.
    
    \item \textbf{Autor Principal:} El nombre del escritor o investigadores primarios.
    
    \item \textbf{Volumen de Distribución:} Cifra numérica que refleja el alcance del libro (unidades vendidas).
    
    \item \textbf{Calificación del Panel de Expertos (Nota de la Crítica):} Una puntuación entera (ej. de 1 a 10) otorgada por un comité evaluador interno, basada en la calidad, originalidad e influencia potencial de la obra.
    
    \item \textbf{Disciplina Principal (Especialidad):} El género al que pertenece la obra (ej: ``Misterio'', ``Thriller'', ``Clásico'', ``Ciencia ficción'', ``Fantasía'', ``Clásico'', ``Realismo mágico'', etc.).
    
    \item \textbf{Extracto (Comentarios):} Nota relevante seleccionada por el equipo de curaduría.
\end{itemize}

Se ha puesto especial énfasis en la eficiencia de ciertas operaciones. Además de la
funcionalidad básica de ``Incorporar Nueva Publicación'' al sistema (la operación ``crear''),
se deben desarrollar las siguientes herramientas de consulta avanzada:

\begin{enumerate}
    \item \textbf{``Recomendación Premium por Autor'' (Compra Segura):} Dada la base de datos
    de publicaciones y el nombre de un Autor Principal, el sistema debe identificar y
    devolver, de forma inmediata, la publicación de dicho autor que ostente la más alta
    Calificación del Panel de Expertos. Esta se considera la ``apuesta segura'' o la obra
    más laureada del autor dentro de la colección.
    
    \item \textbf{``Ranking de Influencia por Disciplina'' (Listado):} Dada la base de datos y una
    Disciplina Principal específica, el sistema debe generar un listado de todas las
    publicaciones pertenecientes a esa disciplina. Dicho listado debe presentarse
    ordenado de mayor a menor según su Volumen de Distribución, facilitando así la
    identificación de las obras más leídas o difundidas.
\end{enumerate}

Un requisito fundamental es que las operaciones de ``Recomendación Premium por
Autor'' y ``Ranking de Influencia por Disciplina'' sean excepcionalmente rápidas
(idealmente, con una complejidad algorítmica constante para su caso promedio, $O(1)$). Esto
es vital para garantizar una experiencia de usuario fluida y para que el sistema pueda ser
utilizado en presentaciones en tiempo real o en portales web de alto tráfico sin demoras
perceptibles.
\end{document}